\section{石碑总表：数学侧 $\leftrightarrow$ 物理侧}
\label{app:stone}

本表是全书词条的压缩版：每行给「数学侧的一个精确对象」、对译机制、以及「物理侧的一个精确对象」。
读法：先看中列，它说明这一步靠什么机制翻译；若中列写「argued」，表示 KW06 只给了物理推导。
页码为原文（hep-th/0604151v3）印刷页，节号与中译本一致。

\begingroup
\small
\begin{longtable}{@{}p{5.4cm}p{4.2cm}p{5.4cm}@{}}
\toprule
\textbf{数学侧} & \textbf{对译机制} & \textbf{物理侧} \\
\midrule
\endhead

\multicolumn{3}{@{}l}{\textit{入口：对偶与参数}}\\
\midrule
$SL(2,\mathbb Z)$ 在 $\tau\in\mathbb H$ 上的作用；$\tau\mapsto-1/(n_{\mathfrak g}\tau)$
& Montonen--Olive（§2.2）&
四维 $\mathcal N=4$ 理论的 $S$-对偶：强弱耦合互换，$g^{2}\leftrightarrow 4\pi/g^{2}$ \\
${}^{L}G$（Langlands 对偶群）：根数据权/余权互换
& 对偶把 $G$ 送到 ${}^{L}G$（§2.2）&
$S$-对偶后的规范群；$\pi_1(G)\leftrightarrow Z({}^{L}G)$ 的 GNO 对偶（§7） \\
典范参数 $\Psi=\Psi(t,\tau)$，$t=v/u\in\mathbb P^1$
& GL 扭化：只依赖 $\Psi$（§3.3--3.5）&
一整套四维拓扑场论 $T(\Psi)$；$t=\pm i$ 复化、$t=\pm1$ 加乘、$t=0,\infty$ 退化 \\
\midrule
\multicolumn{3}{@{}l}{\textit{舞台：Hitchin 模空间}}\\
\midrule
Hitchin 方程 $F-\varphi\wedge\varphi=0,\ D\ast\varphi=0$ 的解模去规范群：$\mathcal M_H(G,C)$
& 紧化到 $C$ 上的 2d TFT（§4.1）&
二维 sigma 模型的目标空间；超凯勒结构的 $I$（Higgs 丛）、$J$（平坦丛）、$K=IJ$ \\
Hitchin 映射 $\pi:\mathcal M_H\to\mathcal B$ 的一般纤维：复环面，复维 $(g-1)\dim G$
& 可积系统 / 自由 $U(1)$ 场（§4.3）&
同一 Hitchin 纤维上的 $T$-对偶：$\mathcal M_H(G)$ 与 $\mathcal M_H({}^{L}G)$ 的纤维互为对偶环面（§5.5） \\
镜像：$A(\mathcal M_H({}^{L}G),\omega_K)\cong B(\mathcal M_H(G),J)$
& 镜像对称（§5.5--5.6）&
$S$-对偶 $\Psi\mapsto-1/(n_{\mathfrak g}\Psi)$ 在二维的化身：$A$-模型与 $B$-模型互换 \\
\midrule
\multicolumn{3}{@{}l}{\textit{膜：边界条件与层}}\\
\midrule
$B$-模型的开弦边界条件；对象级 = 相干层 $F$，态射级 = $\mathrm{Ext}^\bullet$
& B-膜 $\leftrightarrow$ 相干层（§5.5）&
$(B,B,B)$ 型边界条件（$I,J,K$ 中皆为 B-膜）；零膜 $\mathcal O_E$ 支撑在一点上 \\
$A$-模型的开弦边界条件；Lagrange 子流形 + 局部系统
& A-膜 $\leftrightarrow$ Lagrange 膜（§5.5）&
$(A,B,A)$ 或 $(A,A,B)$ 型边界条件（在 $J$ 或 $K$ 中为 A-膜） \\
余各向同性子流形 $Y\subset X$：$TY^{\perp_\omega}\subset TY$
& 余各向同性膜（§11.1 的更精细例子）&
$(A,B,A)$ 或 $(A,A,B)$ 型中支撑高于中间维数的情形；$B_{cc}$ 的支撑对象 \\
支撑在 Hitchin 纤维 $\mathcal F$ 上的秩 1 平坦丛
& type $F$ 膜（磁特征膜）（§9.1）&
$(B,A,A)$ 型边界条件（$I$ 中 B-膜、$J,K$ 中 A-膜）；零膜的 $S$-对偶 \\
\midrule
\multicolumn{3}{@{}l}{\textit{线算子与 Hecke}}\\
\midrule
Hecke 对应 $\mathrm{Bun}_G\xleftarrow{p}\mathcal H_z\xrightarrow{q}\mathrm{Bun}_G$；Hecke 算子 $H_R$
& 线算子 = Hecke 算子（§9.1）&
't Hooft 线算子 $T(U_{z_0})$ 在 $\mathrm{Bun}_G$ 上的作用：在 $z_0$ 处修正丛（§9.1） \\
Hecke 特征层 $\mathcal F_E$：$H_R(\mathcal F_E)\cong V_R\otimes\mathcal F_E$
& 特征膜 = 特征层（§9.1、§11）&
电特征膜：$\mathcal M_H(G)$ 上的膜，被 Wilson 线 $W_R$ 作用成标量（§7--§8） \\
修正空间 $\bar{\mathcal Y}({}^{L}w;z_{0})$ $\cong$ 荷 $k$ 单极子模空间
& 单极子 = Hecke 修正（§10.5，argued 等价）&
单极子冒泡：'t Hooft 线的算子乘积展开（OPE）由单极子模空间的相交上同调给出（§10.4，argued） \\
Nahm 数据：$T_1,T_2,T_3$ 满足 Nahm 方程，极点在区间两端
& ADHMN 构造（定理）&
Bogomolny 方程的解（单极子）；荷 = 留数表示的维数 \\
\midrule
\multicolumn{3}{@{}l}{\textit{终点：几何 Langlands}}\\
\midrule
$D\text{-mod}(\mathrm{Bun}_G)$（扭化 $D$-模）
& 翻译器 $B_{cc}$：$A\mapsto\operatorname{Hom}(B_{cc},A)$（§11.2--11.3，argued）&
$B_{cc}$ 上的开弦：$A$-膜范畴 $\to$ 扭化 $D$-模范畴 \\
$T^{*}\mathrm{Bun}_G$ 上的相干层（B-膜）
& 余切丛的量子化（§11.1）&
$B_{cc}$ 的邻域：$D$-模 = 余切丛非交换形变上的 B-膜 \\
$\mathrm{Loc}_{{}^{L}G}$ 上的准相干层
& 几何 Langlands 对应（§3.1、§11.3）&
电磁对偶后的真空模空间：电特征膜的参数空间 \\
\bottomrule
\end{longtable}
\endgroup

\medskip
\noindent{\small\color{gcol}
注：本表 20 行中，标「定理」或「已证」的 4 行（ADHMN 构造、Hitchin 可积系统、GNO 对偶、
$G=GL_1$ 的类域论对应）；标「argued」的 6 行全部来自 KW06 的物理推导（§9.1、§10.4、§10.5、§11.2--11.3、§12）；
其余为原文明确陈述的对应关系。}
